\section{Introduction}
\label{sec:introduction}

Machine learning has been central to high-energy physics for decades, enabling
reconstruction algorithms, event classification, and detector optimization across
collider experiments. A new frontier has emerged with large language models (LLMs)
developed by major AI laboratories---including
Anthropic~\cite{Anthropic:2025accelerating}, OpenAI~\cite{OpenAI:2025gpt5science},
Google DeepMind~\cite{GoogleDeepMind:2025geminiscience}, and others---that can now
assist directly with scientific reasoning, mathematical proof, and code generation.
A striking asymmetry persists in fundamental physics: nearly all current
AI-assisted scientific workflows operate on the theoretical or computational side.
Recent examples include AI-assisted derivations of new results in
theoretical particle physics~\cite{OpenAI:2025newresult}, design of novel protein
structures for biology and health research~\cite{GoogleDeepMind:2025alphaproteo},
and assistance with complex mathematical simplifications in high-energy physics
calculations. The experimental side of the AI-assisted scientific loop remains
largely absent. To realize a future where AI rapidly accelerates the discovery of
new physics, AI-assisted scientific workflows must be developed and validated
across \textit{both} precision theory \textit{and} precision experiment.

A logical first step toward closing this gap is an agentic analysis of real
experimental data with expert physicist oversight. This is precisely the approach
pursued here: the physicist opened Anthropic Claude Code, brainstormed the
analysis goals, then iteratively prompted the AI agent to execute the full
analysis workflow. No analysis code was written by hand; the physicist's role is
analogous to advising a graduate student, providing expert direction, reviewing
results, and iterating. The agent was supplied with prior ALEPH analysis code,
analysis notes, and the published paper as context. The rationale follows from
the established understanding in the AI literature that few-shot examples and
supplied context improve AI-assisted outputs~\cite{Brown:2020fewshot,Wei:2022cot}.
Pursuing this theory--experimental interface for physics will inevitably deepen
understanding of AI capabilities and their potential to improve AI systems for
science domains where precision matters.

This vision parallels the evolution of AI in neighboring scientific domains. In
materials science and structural biology, AI work has progressed from computational
prediction with post-hoc human validation toward AI-designed experiments executed
by autonomous robotic systems, creating closed feedback loops where AI proposes,
executes, measures, and learns. High-energy physics faces a parallel opportunity.
While we cannot deploy robotic particle detectors (yet), the emergence of open,
curated data, partially enabled by the e+e$^-$ alliance~\cite{EEAlliance:website},
from past experiments---particularly the precision LEP dataset---can serve as a
``pseudo-laboratory'' for developing and validating AI-assisted scientific workflows
that operate across the precision theory--experiment
interface~\cite{Badea:2019vey,Chen:2021uws}. This requires experimental infrastructure
on both sides: precision theoretical predictions and precision experimental measurements
with rigorous covariance quantification, where AI-assisted reasoning is tested and held
accountable to measurable constraints~\cite{Wang:2023aai,Messeri:2024aiu,Woodruff:2026practicalsi}.

Precision QCD provides an excellent domain for this program, precisely because of the
combination of high-precision $e^{+}e^{-}$ data and sophisticated high-precision theory
calculations. Achieving Standard Model predictions at the precision required for future
collider programs, including the High-Luminosity LHC, demands QCD calculations with
uncertainty control matched to percent-level experimental precision. In this regime,
precision experimental measurement becomes validation infrastructure for AI models.
When percent-level uncertainties are coupled to rigorous covariance quantification and
end-to-end reproducibility, high-precision measurements define the testing ground where
AI predictions are made quantitative, falsifiable, and subject to rigorous scientific
accountability. Although there is currently no active high-energy lepton collider,
archived LEP data now provide an unprecedented open-data opportunity for precision
event-shape studies with explicit uncertainty control~\cite{EEAlliance:website,Badea:2019vey,Chen:2021uws}.
Relative to hadron- and nuclear-collision open-data settings, the LEP $e^{+}e^{-}$
environment provides substantially tighter control of soft-event complexity and
initial-state structure, making it a uniquely clean benchmark for both precision QCD
and precision AI model development.

Among event-shape observables, thrust is an ideal first variable for this program.
Thrust $T$ probes both perturbative radiation and non-perturbative dynamics in a
single distribution, is computable at high perturbative order, and has well-established
connections to resummed calculations, making comparison to theory controlled and
interpretable~\cite{PhysRevLett.39.1587,Abbate:2010xh,Heister:2003aj}. It is defined as
$T = \max_{\hat{n}} \left(\sum_i |\vec{p}_i \cdot \hat{n}| / \sum_i |\vec{p}_i|\right)$,
where the sum runs over all selected particles and the thrust axis $\hat{n}$ maximizes
the quantity.

Event-shape observables measured at LEP continue to anchor comparisons between
perturbative calculations, generator modeling, and non-perturbative hadronization
constraints. Recent thrust-based $\alpha_s$ extractions report
$\alpha_s(m_Z)\simeq 0.1136\pm 0.0012$, below the 2023 world-average value
$\alpha_s(m_Z)=0.1180\pm 0.0009$~\cite{Benitez:2024nav,ParticleDataGroup:2022pth}.
This persistent tension, often referred to as the $e^{+}e^{-}$ $\alpha_s$ puzzle,
sharpens the need for renewed measurements with explicit uncertainty control and
publishable covariance information. The thrust peak and logarithmic moments are also
increasingly important for constraining non-perturbative shape-function
effects~\cite{Abbate:2010xh,Assi:2025ibi}.

The immediate objective of this note is to establish a robust detector-corrected
thrust baseline using Iterative Bayesian Unfolding (IBU) implemented in RooUnfold,
with complete validation and systematic propagation. The agent was asked to execute
the full analysis loop: data and simulation processing, unfolding orchestration,
systematic propagation, figure production, and note maintenance with traceable links
to code and outputs. This analysis is also part of a broader program of
AI-assisted scientific workflows in high-energy physics, where agent-oriented
frameworks are developing rapidly alongside AI-assisted theoretical
calculations~\cite{Abolhasani:2023sdl,MacLeod:2020sdl,Chen:2024scienceagentbench,Lu:2024aiscientist,Woodruff:2026practicalsi,Bubeck:2025coscientist,OpenAITheory:2026singleminus,Schwartz:2026cparameterllm}.
In addition to the unfolded spectrum, complete covariance and correlation matrices
are produced for downstream phenomenology use. In the next stage, the same framework
will be extended to additional event-shape observables (C-parameter and heavy-jet
mass), and the approach of using AI-assisted scientific workflows for both theoretical
and experimental analysis will be pursued in a loop between them. Since thrust is
computed from selected charged and neutral particles, the note also includes detailed
\texttt{pwflag}-resolved kinematic comparisons between data and simulation before the
unfolding stage.

The remainder of the document is structured as follows: Sec.~\ref{sec:ai_agent}
defines the AI-agent operating model used in this internal analysis cycle;
Sec.~\ref{sec:aleph} summarizes detector elements most relevant for event-shape
reconstruction; Sec.~\ref{sec:datasets} describes data and simulation samples;
Sec.~\ref{sec:event_selection} details reconstruction and selections, including
particle-category and event-level validation; Sec.~\ref{sec:method} presents
thrust construction and the executable unfolding workflow;
Sec.~\ref{sec:systematics} and Sec.~\ref{sec:results} document uncertainty
propagation and the final thrust baseline outputs; Sec.~\ref{sec:challenges}
closes the note with practical collaboration protocol, persistent analysis
challenges, and mitigation tracking.
